\documentclass{article}
\usepackage[T1]{fontenc}
\usepackage{lmodern}
\usepackage{graphicx}
\usepackage{amsmath,amssymb}
\usepackage[a4paper,margin=2cm]{geometry}
\usepackage[absolute,overlay]{textpos}
\setlength{\TPHorizModule}{1mm}
\setlength{\TPVertModule}{1mm}
\usepackage[indonesian]{babel}
\usepackage{hyperref}
\setlength{\emergencystretch}{2em}
% unit: Halaman 17

\begin{document}

% === Halaman 17 (python/native) ===

• Pembuktian 

            “kunci OTP (acak) + plainteks (tidak  acak) = cipherteks (acak)”

   sebagai berikut:

    Misalkan: 

              C = (P + K) mod 26

  dengan K dipilih secara uniform:  

              P(K = k) = 

1

26​,     k =0,1, …,25.

   Untuk setiap nilai plainteks tertentu P = p, pemetaan

              C = (P + K) mod 26

   merupakan pemetaan satu-ke-satu dari semua kemungkinan K ke semua 

kemungkinan 𝐶.

   Akibatnya:

           P(C = c ∣ P = p) = 

1

26

  untuk setiap c  (yaitu: P(C = ‘a’) = P(C = ‘b’) = ⋯ = P(C = ‘z’) =

1

26

17

\end{document}
